\chapter{Psychiatric and Mental Health Disorders}
\label{ch:psychiatric}
Psychiatric and mental health disorders encompass a broad range of conditions affecting mood, cognition, behavior, and perception. They are among the leading causes of disability worldwide, with major depressive disorder alone affecting over 280 million people globally. The etiology of psychiatric disorders involves complex interactions between genetic predisposition, neurobiological factors, psychological influences, and environmental stressors. Diagnosis follows criteria from the Diagnostic and Statistical Manual of Mental Disorders, Fifth Edition, Text Revision (DSM-5-TR) or the International Classification of Diseases, 11th Revision (ICD-11).

%-------------------------------------------------------------
\section{Anxiety Disorders}
%-------------------------------------------------------------

%-------------------------------------------------------------

%-------------------------------------------------------------

%-------------------------------------------------------------

Anxiety disorders represent a prevalent class of mental health conditions characterized by excessive, persistent fear, anxiety, and behavioral avoidance. Dysregulation within cortico-striato-thalamo-cortical loops and limbic networks impairs emotion regulation and stress coping mechanisms. Effective management integrates evidence-based psychotherapies (cognitive behavioral therapy) and selective serotonin reuptake inhibitors or SNRIs.

%-------------------------------------------------------------

%-------------------------------------------------------------

%-------------------------------------------------------------

%-------------------------------------------------------------

\label{sec:Anxiety Disorders}
%-------------------------------------------------------------%-------------------------------------------------------------

\subsection{Agoraphobia}
\label{sec:Agoraphobia}
Agoraphobia has a lifetime prevalence of 1--3\%, with a mean age of onset in the late 20s. The condition results from conditioned avoidance behaviors and interoceptive hypersensitivity, with pathophysiology overlapping with panic disorder. Clinical features include marked fear or anxiety about being in situations from which escape might be difficult or help might not be available in the event of panic-like symptoms, with typical feared situations including using public transportation, being in open spaces, being in enclosed spaces, standing in line or being in a crowd, or being outside of the home alone. Diagnosis is by clinical interview. Management includes CBT with in vivo exposure and SSRIs/SNRIs. Prognosis is good with treatment.

\subsection{Generalized Anxiety Disorder (GAD)}
\label{sec:Generalized Anxiety Disorder (GAD)}
Generalized anxiety disorder has a lifetime prevalence of 5--9\% and is twice as common in women. The condition results from hyperactivation of the amygdala, prefrontal cortex dysfunction, reduced GABAergic tone, elevated norepinephrine activity, and serotonergic dysregulation. Clinical features include excessive, uncontrollable worry about multiple domains for at least 6 months, occurring more days than not, with associated symptoms including restlessness, easy fatiguability, difficulty concentrating, irritability, muscle tension, and sleep disturbance, with at least three of these required for diagnosis. Diagnosis is by clinical interview; the GAD-7 scale is a validated screening tool. Management includes pharmacotherapy (SSRIs and SNRIs are first-line, with buspirone as an alternative) and psychotherapy (CBT is first-line). Prognosis is variable; many patients respond to treatment but symptoms may be chronic.

\subsection{Panic Disorder / Panic Attacks}
\label{sec:Panic Disorder / Panic Attacks}
Panic disorder has a lifetime prevalence of 3--5\%, with a mean age of onset in early 20s to 30s. The condition results from a hypersensitive fear network centered on the amygdala, locus coeruleus (noradrenergic dysregulation), and periaqueductal gray, with faulty interoceptive conditioning and catastrophic misinterpretation of somatic sensations. Clinical features include recurrent unexpected panic attacks (discrete periods of intense fear or discomfort that reach a peak within minutes, with at least four of: palpitations, sweating, trembling, shortness of breath, sensations of choking, chest pain, nausea, dizziness, chills or heat sensations, paresthesias, derealization, fear of losing control or dying) plus at least one month of persistent concern about additional attacks. Diagnosis is by clinical interview, ruling out cardiac, pulmonary, and endocrine causes. Management includes pharmacotherapy (SSRIs and SNRIs are first-line) and psychotherapy (CBT with interoceptive exposure is highly effective). Prognosis is good with appropriate treatment.

\subsection{Social Anxiety Disorder (Social Phobia)}
\label{sec:Social Anxiety Disorder (Social Phobia)}
Social anxiety disorder has a lifetime prevalence of 7--13\%, making it one of the most common anxiety disorders, typically beginning in adolescence with peak onset at age 13. The condition results from amygdala hyperreactivity, insula hyperactivity, and altered default-mode network connectivity. Clinical features include marked fear or anxiety about one or more social situations in which the individual is exposed to possible scrutiny by others, with fear of being negatively evaluated, embarrassment, rejection, or offending others, with the fear being out of proportion to the actual threat and persistent (at least 6 months). Physical symptoms (blushing, trembling, sweating, palpitations) are common. Diagnosis is by clinical interview. Management includes SSRIs and SNRIs as first-line pharmacotherapy, and CBT (cognitive restructuring and exposure therapy) as first-line psychotherapy. Prognosis is good with treatment.

\subsection{Specific Phobias}
\label{sec:Specific Phobias}
Specific phobias have a lifetime prevalence of 10--13\%, making them the most common anxiety disorder in the general population. The condition results from conditioned fear responses. Clinical features include marked, persistent (at least 6 months) fear or anxiety about a specific object or situation (e.g., flying, heights, animals, receiving an injection, blood, enclosed spaces), with the phobic stimulus almost invariably provoking an immediate fear response out of proportion to the actual danger. Common subtypes include animal type, natural environment type, blood-injection-injury type (unique vasovagal response with fainting in up to 75\%), situational type, and other type. Diagnosis is by clinical interview. Management of choice is CBT with exposure therapy (systematic desensitization). Prognosis is excellent with treatment.

%-------------------------------------------------------------
\section{Disruptive, Impulse-Control, and Conduct Disorders}
%-------------------------------------------------------------

%-------------------------------------------------------------

%-------------------------------------------------------------

%-------------------------------------------------------------

%-------------------------------------------------------------

%-------------------------------------------------------------

%-------------------------------------------------------------

%-------------------------------------------------------------

\label{sec:Disruptive, Impulse-Control, and Conduct Disorders}
%-------------------------------------------------------------%-------------------------------------------------------------

Disruptive, impulse-control, and conduct disorders include oppositional defiant disorder (ODD), conduct disorder (CD), intermittent explosive disorder (IED), pyromania, and kleptomania.

\textbf{Oppositional Defiant Disorder (ODD):} Prevalence 2--11\%, more common in males before adolescence. Characterized by a pattern of angry/irritable mood, argumentative/defiant behavior, and vindictiveness lasting at least 6 months in children and adolescents. Often a precursor to conduct disorder. Treatment: parent management training, collaborative problem solving, CBT for the child.

\textbf{Conduct Disorder (CD):} Prevalence 2--10\%, male predominance. Characterized by a repetitive and persistent pattern of behavior violating the rights of others or major age-appropriate societal norms---aggression to people and animals, destruction of property, deceitfulness or theft, and serious rule violations. Onset type: childhood-onset (before age 10, more severe, more likely to develop antisocial personality disorder) and adolescent-onset (better prognosis, limited to adolescence). Callous-unemotional traits specify a more severe, treatment-refractory subtype. Treatment: multisystemic therapy (MST), functional family therapy, parent management training---the Oregon model. Pharmacotherapy targets aggression (risperidone, methylphenidate for comorbid ADHD).

\textbf{Intermittent Explosive Disorder (IED):} Prevalence 2--3\%. Recurrent behavioral outbursts representing a failure to control aggressive impulses, resulting in verbal aggression or physical aggression against property, animals, or individuals. The outbursts are grossly out of proportion to the provocation. Treatment: CBT (anger management, cognitive restructuring), SSRIs (fluoxetine), mood stabilizers, and atypical antipsychotics.

\textbf{Kleptomania:} Recurrent failure to resist impulses to steal items not needed for personal use or monetary value, with increasing tension before the theft and pleasure, gratification, or relief during it.

\textbf{Pyromania:} Deliberate and purposeful fire-setting on more than one occasion, with tension or affective arousal before the act and fascination with, interest in, or attraction to fire and its situational contexts.

Both kleptomania and pyromania are rare and typically require specialized CBT (Cue Exposure and Response Prevention---CERP for kleptomania) and SSRIs.

%-------------------------------------------------------------
\section{Dissociative Disorders}
%-------------------------------------------------------------

%-------------------------------------------------------------

%-------------------------------------------------------------

%-------------------------------------------------------------

%-------------------------------------------------------------

%-------------------------------------------------------------

%-------------------------------------------------------------

%-------------------------------------------------------------

\label{sec:Dissociative Disorders}
%-------------------------------------------------------------%-------------------------------------------------------------

\subsection{Dissociative Amnesia}
\label{sec:Dissociative Amnesia}
Dissociative amnesia has a prevalence of 1--3\% in the general population. The condition results from stress-induced disruption of hippocampal and prefrontal memory encoding and retrieval processes, typically associated with trauma (sexual assault, combat, natural disaster, interpersonal violence). Clinical features include inability to recall important autobiographical information, usually of a traumatic or stressful nature, that is inconsistent with ordinary forgetting, with the amnesia possibly being localized, selective, generalized, or systematized. Diagnosis is clinical, with neuroimaging and EEG to rule out organic amnesia. Management includes psychotherapy (psychodynamic, CBT, hypnotherapy) focused on recovery from trauma and integration. Prognosis is good for localized and selective amnesia---memory typically returns spontaneously or with treatment.

\subsection{Dissociative Identity Disorder}
\label{sec:Dissociative Identity Disorder}
Dissociative identity disorder has a prevalence of 1--1.5\% in the general population, though it is frequently underdiagnosed. The condition typically arises from severe, repetitive, early childhood trauma before age 6--9, when the child's personality is still developing and dissociation serves as a psychological defense, with evidence of functional differences between identities on neuroimaging. Clinical features include disruption of identity involving two or more distinct personality states (alter identities), each with its own relatively enduring pattern of perceiving, relating to, and thinking about the environment and self, accompanied by recurrent gaps in the recall of everyday events, important personal information, and/or traumatic events. Diagnosis is by clinical interview (SCID-D). Management is long-term, phase-oriented treatment focusing on safety/stabilization, trauma processing, and identity integration. Prognosis is variable; treatment aims for improved cooperation between identities or integration.

%-------------------------------------------------------------
\section{Eating Disorders}
%-------------------------------------------------------------

%-------------------------------------------------------------

%-------------------------------------------------------------

%-------------------------------------------------------------

Eating disorders are complex psychiatric conditions characterized by severe disturbances in eating behaviors, body image perception, and metabolic regulation. Complex interactions between genetic vulnerability, neurochemical alterations, and sociocultural pressures drive persistent caloric restriction or binge-purge cycles. Multidisciplinary management encompasses nutritional rehabilitation, medical stabilization, individual psychotherapy, and family-based treatment.

%-------------------------------------------------------------

%-------------------------------------------------------------

%-------------------------------------------------------------

%-------------------------------------------------------------

\label{sec:Eating Disorders}
%-------------------------------------------------------------%-------------------------------------------------------------

\subsection{Anorexia Nervosa}
\label{sec:Anorexia Nervosa}
Anorexia nervosa has a lifetime prevalence of 0.5--1\% in women and 0.1--0.3\% in men, with a 10:1 female-to-male ratio, and age of onset typically in adolescence (peak 15--19 years). The condition results from serotonergic and dopaminergic dysregulation, aberrant reward processing, and sociocultural factors (thin-ideal internalization). Clinical features include restriction of energy intake relative to requirements leading to significantly low body weight, intense fear of gaining weight or becoming fat (even when underweight), and disturbance in self-perceived weight or shape, with subtypes including restricting type and binge-eating/purging type. Medical complications include bradycardia, hypotension, hypothermia, osteopenia/osteoporosis, lanugo, cardiac arrhythmias (prolonged QT interval), electrolyte disturbances, refeeding syndrome, and structural brain changes. Diagnosis is by clinical interview. Management includes medical stabilization, nutritional rehabilitation, psychotherapy (family-based treatment for adolescents, CBT-E for adults), and olanzapine for weight gain. Prognosis is variable: 50\% recover, 30\% improve, 20\% remain chronic, with mortality of 5--10\%.

\subsection{Binge-Eating Disorder}
\label{sec:Binge-Eating Disorder}
Binge-eating disorder has a lifetime prevalence of 2--3\% in the general population and is the most common eating disorder, more common in women but with a narrower sex ratio than anorexia and bulimia (3:2 female-to-male). The condition results from altered reward sensitivity (dopamine dysfunction in striatum), impulsivity, and emotional dysregulation. Clinical features include recurrent episodes of binge eating (loss-of-control eating of objectively large amounts of food) without regular compensatory behaviors, associated with eating much more rapidly than normal, eating until uncomfortably full, eating large amounts when not physically hungry, eating alone due to embarrassment, and feeling disgusted/depressed/guilty afterward, occurring at least once a week for 3 months. BED is strongly associated with obesity (60--80\% are obese) and metabolic syndrome. Diagnosis is by clinical interview. Management includes CBT-E, IPT, and lisdexamfetamine (FDA-approved for moderate-to-severe BED). Prognosis is better than for anorexia or bulimia; 50--70\% achieve remission with structured treatment.

\subsection{Bulimia Nervosa}
\label{sec:Bulimia Nervosa}
Bulimia nervosa has a lifetime prevalence of 1--1.5\% in women and 0.5\% in men, with an age of onset in late adolescence to early adulthood. The condition results from altered reward circuitry, deficits in inhibitory control, and interoceptive dysfunction. Clinical features include recurrent episodes of binge eating (eating a large amount of food in a discrete period with a sense of loss of control) followed by recurrent inappropriate compensatory behaviors to prevent weight gain (self-induced vomiting, laxative/diuretic misuse, fasting, excessive exercise), occurring at least once a week for 3 months, with self-evaluation unduly influenced by body shape and weight. Medical complications include electrolyte disturbances (hypokalemic hypochloremic metabolic alkalosis), Russell's sign, parotid gland enlargement, dental erosion, Mallory-Weiss tears, esophageal rupture, and ipecac cardiomyopathy. Diagnosis is by clinical interview. Management includes CBT-E (first-line), IPT (second-line), and SSRI pharmacotherapy (high-dose fluoxetine 60 mg/day is FDA-approved). Prognosis: 50--60\% achieve recovery; relapse rates are moderate.

%-------------------------------------------------------------
\section{Mood Disorders}
%-------------------------------------------------------------

%-------------------------------------------------------------

%-------------------------------------------------------------

%-------------------------------------------------------------

Mood disorders, principally major depressive disorder and bipolar spectrum disorders, involve pervasive dysregulation of affect, energy levels, and vegetative functions. Neurobiological alterations in monoaminergic neurotransmission, neuroplasticity, and neuroendocrine axes contribute to persistent emotional suffering and functional disability. Treatment combines evidence-based psychotherapies, mood-stabilizing medications, antidepressants, and neurostimulation therapies when indicated.

%-------------------------------------------------------------

%-------------------------------------------------------------

%-------------------------------------------------------------

%-------------------------------------------------------------

\label{sec:Mood Disorders}
%-------------------------------------------------------------%-------------------------------------------------------------

\subsection{Bipolar I Disorder}
\label{sec:Bipolar I Disorder}
Bipolar I disorder has a lifetime prevalence of approximately 1\% worldwide, with equal distribution across sexes. The condition results from dysregulation of intracellular signaling pathways (GSK-3, PKC, protein kinase A) and circadian rhythm disruption, with strong genetic factors (80--90\% heritability). Clinical features include at least one lifetime manic episode lasting at least 1 week (or any duration if hospitalization is required), characterized by elevated, expansive, or irritable mood plus increased goal-directed activity or energy, with at least three of the following (four if mood is only irritable): inflated self-esteem or grandiosity, decreased need for sleep, pressured speech, flight of ideas or racing thoughts, distractibility, psychomotor agitation, and excessive involvement in high-risk activities. Diagnosis is by clinical interview using DSM-5-TR criteria. Management of acute mania includes a mood stabilizer plus an atypical antipsychotic or benzodiazepine; lithium remains the gold-standard mood stabilizer with demonstrated antisuicidal effects, with other first-line treatments including valproate and atypical antipsychotics (olanzapine, quetiapine, risperidone, aripiprazole). Prognosis is variable, with most patients requiring long-term maintenance therapy.

\subsection{Bipolar II Disorder}
\label{sec:Bipolar II Disorder}
Bipolar II disorder has a lifetime prevalence of approximately 0.5\%, characterized by one or more hypomanic episodes (distinct period of elevated, expansive, or irritable mood lasting at least 4 consecutive days, plus similar symptoms as mania but of lesser severity, without marked functional impairment or psychosis) and one or more major depressive episodes. The condition results from similar pathophysiological mechanisms as Bipolar I. Clinical features include patients spending significantly more time in depressive episodes than hypomanic episodes, often leading to underdiagnosis. Diagnosis requires the absence of full manic episodes, distinguishing it from Bipolar I. Management includes mood stabilizers (lithium, lamotrigine, valproate) and atypical antipsychotics (quetiapine, which is FDA-approved for bipolar II depression), with careful use of antidepressants with mood stabilizer coverage. Prognosis is variable but generally better than Bipolar I.

\subsection{Cyclothymic Disorder}
\label{sec:Cyclothymic Disorder}
Cyclothymic disorder has a lifetime prevalence of 0.5--1\%, characterized by at least 2 years (1 year in children and adolescents) of numerous periods of hypomanic symptoms that do not meet full criteria for a hypomanic episode and numerous periods of depressive symptoms that do not meet full criteria for a major depressive episode, with symptoms present at least half the time. The condition is often described as a temperamental or personality disposition with mood instability. Clinical features include chronic mood instability without meeting full criteria for hypomanic or depressive episodes. Diagnosis is by clinical interview. Management includes mood stabilizers (lithium, lamotrigine) and psychotherapy (IPSRT, CBT). Prognosis is variable, with up to 30\% of patients eventually developing Bipolar I or Bipolar II disorder.

\subsection{Major Depressive Disorder (MDD)}
\label{sec:Major Depressive Disorder (MDD)}
Major depressive disorder is a prevalent psychiatric disorder with a lifetime prevalence of approximately 15--20\% and a 12-month prevalence of 7\%, twice as common in women as in men, with a peak age of onset of 20--40 years. The condition results from dysregulation of monoamine neurotransmitter systems (serotonin, norepinephrine, dopamine), hypothalamic-pituitary-adrenal (HPA) axis hyperactivity with elevated cortisol, reduced hippocampal volume due to chronic stress-induced neurotoxicity, decreased brain-derived neurotrophic factor (BDNF), and altered functional connectivity in frontolimbic circuits, with genetic factors accounting for 37\% of risk. Clinical features include depressed mood (sadness, emptiness, hopelessness) nearly every day, anhedonia (markedly diminished interest or pleasure in activities), significant weight change, insomnia or hypersomnia, psychomotor agitation or retardation, fatigue, feelings of worthlessness or excessive guilt, diminished concentration, and recurrent thoughts of death or suicidal ideation, with DSM-5-TR criteria requiring at least five of these symptoms (including depressed mood or anhedonia) present during the same two-week period. Diagnosis is by clinical interview using DSM-5-TR criteria, with screening tools including the PHQ-9 and BDI-II. Management is multimodal: pharmacotherapy (SSRIs and SNRIs are first-line), psychotherapy (CBT, IPT), somatic therapies (ECT for severe cases, TMS, ketamine/esketamine for treatment-resistant depression), and lifestyle interventions. Prognosis is variable; many patients respond to treatment but recurrence is common.

\subsection{Persistent Depressive Disorder (Dysthymia)}
\label{sec:Persistent Depressive Disorder (Dysthymia)}
Persistent depressive disorder (dysthymia) is a chronic form of depression with symptoms present most of the day, for more days than not, for at least 2 years (1 year in children and adolescents), with a lifetime prevalence of 3--6\%. The condition is characterized by depressed mood plus at least two of the following: poor appetite or overeating, insomnia or hypersomnia, low energy or fatigue, low self-esteem, poor concentration, or feelings of hopelessness, with symptoms being less severe than MDD but more persistent. Clinical features include being ``depressed for as long as I can remember,'' with double depression referring to MDD superimposed on pre-existing dysthymia. Diagnosis is by clinical interview. Management includes SSRIs/SNRIs (same as MDD) and psychotherapy (CBT, IPT). Prognosis is guarded, with high rates of recurrence, and treatment is often recommended for at least 2 years.

%-------------------------------------------------------------
\section{Neurodevelopmental Disorders}
%-------------------------------------------------------------

%-------------------------------------------------------------

%-------------------------------------------------------------

%-------------------------------------------------------------

%-------------------------------------------------------------

%-------------------------------------------------------------

%-------------------------------------------------------------

%-------------------------------------------------------------

\label{sec:Neurodevelopmental Disorders}
%-------------------------------------------------------------%-------------------------------------------------------------

\subsection{Attention-Deficit/Hyperactivity Disorder (ADHD)}
\label{sec:Attention-Deficit/Hyperactivity Disorder (ADHD)}
Attention-deficit/hyperactivity disorder has a worldwide prevalence of 5--7\% in children and 2.5--3\% in adults, more common in boys in childhood (3:1 ratio) but approaching parity in adulthood. The condition results from striatal and prefrontal dopaminergic and noradrenergic dysregulation, with reduced brain volume in prefrontal cortex, caudate, and cerebellum, delayed cortical maturation, and high genetic heritability (75--80\%). Clinical features include a persistent pattern of inattention, hyperactivity, and impulsivity that interferes with functioning, with three presentations: predominantly inattentive, predominantly hyperactive-impulsive, and combined presentation, with onset before age 12. Diagnosis is clinical, based on detailed developmental history, collateral reports, and behavior rating scales. Management includes pharmacotherapy (stimulants as first-line: methylphenidate and amphetamine formulations; non-stimulants as second-line: atomoxetine, guanfacine, clonidine) and psychosocial interventions (behavioral parent training, CBT, executive function coaching). Prognosis: 50--65\% of children continue to have impairing symptoms into adulthood.

\subsection{Autism Spectrum Disorder (ASD)}
\label{sec:Autism Spectrum Disorder (ASD)}
Autism spectrum disorder has a prevalence of 1--2\% worldwide, with a male-to-female ratio of 3--4:1. The condition results from dysregulation of synaptic pruning and neural connectivity, an imbalance of excitation/inhibition (particularly GABAergic and glutamatergic systems), aberrant brain growth, and alterations in frontotemporal, frontoparietal, and interhemispheric connectivity, with strong genetic factors (heritability 80--90\%). Clinical features include persistent deficits in social communication and social interaction across multiple contexts, and restricted, repetitive patterns of behavior, interests, or activities, with three levels of severity requiring support. Social communication deficits include reduced social-emotional reciprocity, deficits in nonverbal communicative behaviors, and deficits in developing relationships. Restricted, repetitive patterns manifest as stereotyped speech/gait, insistence on sameness, highly restricted fixated interests, and hyper- or hyporeactivity to sensory input. Intellectual disability is present in 30--50\%, and language impairment in 25--30\%. Diagnosis is clinical. Management includes early intensive behavioral intervention (EIBI/ABA), social skills training, speech and language therapy, and pharmacotherapy for comorbid symptoms (atypical antipsychotics for irritability/aggression, SSRIs for anxiety/depression). Prognosis is highly variable; early intervention significantly improves functional outcomes.

%-------------------------------------------------------------
\section{Obsessive-Compulsive and Related Disorders}
%-------------------------------------------------------------

%-------------------------------------------------------------

%-------------------------------------------------------------

%-------------------------------------------------------------

%-------------------------------------------------------------

%-------------------------------------------------------------

%-------------------------------------------------------------

%-------------------------------------------------------------

\label{sec:Obsessive-Compulsive and Related Disorders}
%-------------------------------------------------------------%-------------------------------------------------------------

\subsection{Body Dysmorphic Disorder}
\label{sec:Body Dysmorphic Disorder}
Body dysmorphic disorder has a prevalence of 1.7--2.4\% in the general population, with higher rates in dermatology and cosmetic surgery settings (8--15\%). The condition results from distorted body image perception. Clinical features include a preoccupation with one or more perceived defects or flaws in physical appearance that are not observable or appear only slightly noticeable to others, with repetitive behaviors (mirror checking, excessive grooming, skin picking, reassurance seeking) or mental acts (comparing appearance with others) in response to the appearance concerns, with the preoccupation being time-consuming (at least 1 hour per day). Common concerns involve the skin, hair, nose, eyes, teeth, and weight or muscle size (in muscle dysmorphia). Diagnosis is by clinical interview. Management includes SSRIs at high doses (similar to OCD) and CBT tailored for BDD. Prognosis is chronic without treatment; cosmetic procedures generally do not help.

\subsection{Hoarding Disorder}
\label{sec:Hoarding Disorder}
Hoarding disorder has a prevalence of 2--6\% in the general population, increasing with age (particularly in those over 55). Clinical features include persistent difficulty discarding or parting with possessions, regardless of their actual value, due to a perceived need to save them and distress associated with discarding them, leading to accumulation of possessions that congest and clutter active living areas, causing significant distress or functional impairment including health and safety risks. Diagnosis is by clinical interview. Management includes CBT specifically adapted for hoarding (skills training in decision-making, organizing, categorizing; motivational interviewing; in-home decluttering sessions). Prognosis is variable, with high rates of relapse.

\subsection{Obsessive-Compulsive Disorder (OCD)}
\label{sec:Obsessive-Compulsive Disorder (OCD)}
Obsessive-compulsive disorder has a lifetime prevalence of 2--3\% worldwide, with equal distribution across sexes, mean age of onset of 19.5 years, and a bimodal distribution (early-onset peaking in childhood/adolescence, late-onset peaking in the third decade). The condition results from hyperactivity in the cortico-striato-thalamo-cortical (CSTC) loop, specifically the orbitofrontal cortex, anterior cingulate cortex, and caudate nucleus, with glutamatergic hyperactivity and reduced serotonin function. Clinical features include obsessions (recurrent, persistent, intrusive, and unwanted thoughts, urges, or images) and/or compulsions (repetitive behaviors or mental acts that the individual feels driven to perform), with common obsession themes including contamination, harm, symmetry and order, and taboo thoughts, and common compulsions including washing, checking, counting, repeating, and ordering. Diagnosis is by clinical interview. Management includes pharmacotherapy (SSRIs are first-line, typically at higher doses than for depression, with clomipramine as an effective alternative) and psychotherapy (exposure and response prevention [ERP] is the gold-standard psychotherapy). Prognosis is variable: 30--40\% have a chronic course, 40--50\% have moderate improvement, 10--20\% have refractory disease.

\subsection{Tourette Syndrome}
\label{sec:Tourette Syndrome}
Tourette syndrome is a neurodevelopmental disorder characterized by multiple motor tics and one or more vocal tics persisting for more than one year, with onset before age 18, affecting 0.3--0.8\% of children. The condition results from dysfunction of cortico-striato-thalamo-cortical circuits, involving dopaminergic and GABAergic dysregulation, with genetic contributions (inherited as autosomal dominant with variable penetrance; SLITRK1, HDC, CNTN6 genes implicated). Clinical features include motor tics (eye blinking, facial grimacing, head jerking, shoulder shrugging, complex tics such as hopping, touching, or echopraxia) and vocal tics (throat clearing, grunting, sniffing, coprolalia, palilalia, echolalia). Tics wax and wane in severity, are preceded by premonitory urges, and can be temporarily suppressed. Comorbidities include ADHD (50--60\%), OCD (30--60\%), anxiety, and autism spectrum disorder. Diagnosis is clinical, based on DSM-5-TR criteria. Management includes psychoeducation, habit reversal training and comprehensive behavioral intervention for tics (CBIT), pharmacotherapy for moderate to severe tics (alpha-2 agonists: clonidine, guanfacine as first-line), and antipsychotics (aripiprazole, risperidone, haloperidol) for refractory tics. Prognosis is favorable; most patients experience reduction in tic severity by early adulthood.

%-------------------------------------------------------------
\section{Personality Disorders}
%-------------------------------------------------------------

%-------------------------------------------------------------

%-------------------------------------------------------------

%-------------------------------------------------------------

%-------------------------------------------------------------

%-------------------------------------------------------------

%-------------------------------------------------------------

%-------------------------------------------------------------

\label{sec:Personality Disorders}
%-------------------------------------------------------------%-------------------------------------------------------------

\subsection{Antisocial Personality Disorder (ASPD)}
\label{sec:Antisocial Personality Disorder (ASPD)}
Antisocial personality disorder has a prevalence of 0.5--3\% in the general population, with a male-to-female ratio of 3:1, and is higher in prison populations (50--80\%). The condition results from reduced amygdala volume and reactivity, ventromedial prefrontal cortex dysfunction, reduced autonomic responsivity to distress cues, and reduced gray matter volume in the orbitofrontal cortex and anterior insula, with genetic factors being strong (heritability 50--60\%). Clinical features include a pervasive pattern of disregard for and violation of the rights of others, occurring since age 15, with at least three of the following: failure to conform to social norms, deceitfulness, impulsivity, irritability and aggressiveness, reckless disregard for safety, consistent irresponsibility, and lack of remorse, with the individual being at least 18 years old and having evidence of conduct disorder before age 15. Diagnosis is by clinical interview. Management is challenging---no FDA-approved medications for ASPD, with pharmacotherapy targeting aggression and impulsivity (lithium, valproate, atypical antipsychotics). Prognosis is poor; symptoms may attenuate with age but functional impairment persists.

\subsection{Borderline Personality Disorder (BPD)}
\label{sec:Borderline Personality Disorder (BPD)}
Borderline personality disorder has a prevalence of 1--2.5\% in the general population and 10--20\% in outpatient psychiatric settings. The condition results from a hyperreactive amygdala (heightened emotional sensitivity), impaired prefrontal regulation of emotion, reduced hippocampal and amygdala volumes, and altered oxytocin and endogenous opioid systems, with adverse childhood experiences being strongly associated. Clinical features include a pervasive pattern of instability in interpersonal relationships, self-image, affect, and marked impulsivity, with DSM-5-TR criteria requiring at least five of the following: frantic efforts to avoid real or imagined abandonment, unstable and intense interpersonal relationships (alternating between idealization and devaluation---splitting), marked identity disturbance, impulsivity, recurrent suicidal behavior or non-suicidal self-injury (NSSI), affective instability, chronic feelings of emptiness, inappropriate anger, and transient stress-related paranoid ideation or dissociative symptoms. Diagnosis is by clinical interview. Management includes dialectical behavior therapy (DBT) as first-line psychotherapy, mentalization-based treatment (MBT), and symptom-targeted pharmacotherapy (SSRIs, mood stabilizers, low-dose atypical antipsychotics). Prognosis: 80\% achieve symptomatic remission over 10-year follow-up; suicide risk is 5--10\%.

\subsection{Narcissistic Personality Disorder}
\label{sec:Narcissistic Personality Disorder}
Narcissistic personality disorder has a prevalence of 0.5--1\% in the general population, with a male predominance (60--75\%). The condition may involve exaggerated self-enhancement to compensate for underlying fragile self-esteem, deficits in mentalization, and altered frontolimbic connectivity. Clinical features include a pervasive pattern of grandiosity (in fantasy or behavior), need for admiration, and lack of empathy, with required criteria including at least five of the following: grandiose sense of self-importance, preoccupation with fantasies of unlimited success, belief of being special and unique, need for excessive admiration, sense of entitlement, interpersonal exploitativeness, lack of empathy, envy of others, and arrogant behaviors or attitudes. Vulnerable (covert) narcissism is a subtype characterized by hypersensitivity to evaluation and chronic feelings of shame. Diagnosis is by clinical interview. Management includes psychotherapy (transfer-focused, mentalization-based, schema therapy). Prognosis is variable; grandiosity may diminish with age.

%-------------------------------------------------------------
\section{Schizophrenia Spectrum and Other Psychotic Disorders}
%-------------------------------------------------------------

%-------------------------------------------------------------

%-------------------------------------------------------------

%-------------------------------------------------------------

%-------------------------------------------------------------

%-------------------------------------------------------------

%-------------------------------------------------------------

%-------------------------------------------------------------

\label{sec:Schizophrenia Spectrum and Other Psychotic Disorders}
%-------------------------------------------------------------%-------------------------------------------------------------

\subsection{Delusional Disorder}
\label{sec:Delusional Disorder}
Delusional disorder has a lifetime prevalence of 0.2\%, with a later age of onset (typically 40--60 years) compared to schizophrenia. Clinical features include one or more delusions that persist for at least 1 month, without meeting full criteria for schizophrenia (hallucinations, if present, are not prominent and are related to the delusional theme; negative symptoms are absent; functioning is not markedly impaired). Common subtypes include persecutory (most common), jealous, erotomanic (de Clérambault syndrome), grandiose, and somatic. Diagnosis is by clinical interview. Management includes antipsychotics (low-dose). Prognosis is better overall than schizophrenia.

\subsection{Schizoaffective Disorder}
\label{sec:Schizoaffective Disorder}
Schizoaffective disorder has a lifetime prevalence of 0.3\%. The condition is characterized by an uninterrupted period of illness with concurrent features of schizophrenia (active-phase symptoms of delusions, hallucinations, disorganized speech, negative symptoms) and a major mood episode (depressive or manic) that occurs for the majority of the total duration of the illness. Clinical features require that delusions or hallucinations occur without prominent mood symptoms for at least 2 weeks. Diagnosis is by clinical interview. Management combines antipsychotics and mood stabilizers/antidepressants guided by the predominant mood polarity. Prognosis is intermediate between schizophrenia and mood disorders.

\subsection{Schizophrenia}
\label{sec:Schizophrenia}
Schizophrenia has a lifetime prevalence of 0.3--0.7\%, with a male-to-female ratio of 1.4:1, with men typically presenting earlier (peak onset 20--28 years) than women (peak onset 25--35 years). The condition results from dopaminergic hyperactivity in the mesolimbic pathway (positive symptoms), dopaminergic hypoactivity in the mesocortical pathway (negative symptoms and cognitive deficits), glutamatergic NMDA receptor hypofunction, and widespread gray matter loss, with genetic factors accounting for 80\% of risk. Clinical features include two or more of the following for a significant portion of time during a 1-month period (at least one must be delusions, hallucinations, or disorganized speech): delusions (fixed false beliefs, often bizarre), hallucinations (auditory most common), disorganized speech (loose associations, tangentiality), grossly disorganized or catatonic behavior, and negative symptoms (blunted affect, avolition, alogia, anhedonia, asociality), with continuous signs persisting for at least 6 months. Diagnosis is clinical---no biomarkers are available. Management includes antipsychotics (both first-generation and second-generation, with clozapine being the only FDA-approved medication for treatment-resistant schizophrenia), cognitive-behavioral therapy for psychosis (CBTp), psychosocial interventions, and cognitive remediation therapy. Prognosis is variable: 20--30\% have a good outcome, with a chronic course and functional decline in many, and suicide risk of 5--10\%.

%-------------------------------------------------------------
\section{Sleep-Wake Disorders}
%-------------------------------------------------------------

%-------------------------------------------------------------

%-------------------------------------------------------------

%-------------------------------------------------------------

%-------------------------------------------------------------

%-------------------------------------------------------------

%-------------------------------------------------------------

%-------------------------------------------------------------

\label{sec:Sleep-Wake Disorders}
%-------------------------------------------------------------%-------------------------------------------------------------

Sleep-wake disorders are covered in detail in the sleep disorders chapter (Ch.~\ref{ch:sleep}). Disorders that frequently co-occur with psychiatric conditions include insomnia disorder (50--80\% prevalence in depression, 70--90\% in anxiety disorders), hypersomnolence disorder (comorbid with mood disorders), circadian rhythm sleep-wake disorders (frequent in bipolar disorder), parasomnias (associated with PTSD), and nightmare disorder (hallmark of PTSD, also associated with borderline personality disorder). Screening for sleep disorders is essential in all psychiatric evaluations. CBT for insomnia (CBT-I) is the first-line treatment for chronic insomnia, including in psychiatric populations.

%-------------------------------------------------------------
\section{Somatic Symptom and Related Disorders}
%-------------------------------------------------------------

%-------------------------------------------------------------

%-------------------------------------------------------------

%-------------------------------------------------------------

%-------------------------------------------------------------

%-------------------------------------------------------------

%-------------------------------------------------------------

%-------------------------------------------------------------

\label{sec:Somatic Symptom and Related Disorders}
%-------------------------------------------------------------%-------------------------------------------------------------

\subsection{Illness Anxiety Disorder (Hypochondriasis)}
\label{sec:Illness Anxiety Disorder (Hypochondriasis)}
Illness anxiety disorder has a prevalence of 1--5\% in the general population, with equal distribution across sexes. The condition results from interoceptive amplification, catastrophic misinterpretation, and memory biases for threat-related health information. Clinical features include preoccupation with having or acquiring a serious illness, with somatic symptoms being absent or mild, and high levels of anxiety about health, with the care-seeking type seeking repeated medical evaluations and the care-avoidant type avoiding medical care due to fear of discovering serious illness. Diagnosis is by clinical interview. Management includes CBT with psychoeducation, cognitive restructuring, behavioral experiments, and response prevention. Prognosis: chronic fluctuating course.

\subsection{Somatic Symptom Disorder}
\label{sec:Somatic Symptom Disorder}
Somatic symptom disorder has an estimated prevalence of 5--7\% in the general population and is more common in women. The condition results from hypervigilance to bodily sensations, amplification of normal physiological signals, maladaptive health-related cognitions, and altered brain connectivity in insula, anterior cingulate, and prefrontal cortex. Clinical features include one or more somatic symptoms that are distressing or result in significant disruption of daily life, with excessive thoughts, feelings, or behaviors related to the somatic symptoms manifested by disproportionate and persistent thoughts about the seriousness of symptoms, persistently high level of anxiety about health, or excessive time and energy devoted to these symptoms. Diagnosis is by clinical interview. Management includes CBT (cognitive restructuring, attention retraining, behavioral activation) and pharmacotherapy focusing on comorbid depression and anxiety (SSRIs). Prognosis is chronic but function can improve with treatment.

%-------------------------------------------------------------
\section{Substance-Related and Addictive Disorders}
%-------------------------------------------------------------

%-------------------------------------------------------------

%-------------------------------------------------------------

%-------------------------------------------------------------

%-------------------------------------------------------------

%-------------------------------------------------------------

%-------------------------------------------------------------

%-------------------------------------------------------------

\label{sec:Substance-Related and Addictive Disorders}
%-------------------------------------------------------------%-------------------------------------------------------------

\subsection{Alcohol Use Disorder}
\label{sec:Alcohol Use Disorder}
Alcohol use disorder is one of the most prevalent psychiatric disorders, with a 12-month prevalence of 5--8\% and lifetime prevalence of 20--30\% in the United States, more common in men but increasing in women. The condition results from dysregulation of the brain reward system (mesolimbic dopamine pathway), with chronic alcohol use producing tolerance, withdrawal (GABA receptor downregulation and NMDA receptor upregulation), and sensitization of stress systems, with genetic factors contributing 50--60\% of risk. Clinical features include at least 2 of 11 DSM-5-TR criteria occurring within a 12-month period, including alcohol taken in larger amounts or over longer periods than intended, persistent desire to cut down, craving, recurrent use resulting in failure to fulfill obligations, continued use despite problems, tolerance, and withdrawal. Medical complications include alcoholic hepatitis, cirrhosis, pancreatitis, Wernicke encephalopathy (confusion, ataxia, ophthalmoplegia---treated with IV thiamine), Korsakoff syndrome, alcoholic cardiomyopathy, peripheral neuropathy, and increased cancer risk. Diagnosis is by clinical interview. Management includes detoxification (benzodiazepines for withdrawal prevention), pharmacotherapy for relapse prevention (naltrexone, acamprosate, disulfiram), and psychosocial interventions (CBT, MET, 12-step facilitation). Prognosis: chronic relapsing condition; 30--50\% achieve sustained sobriety with treatment.

\subsection{Opioid Use Disorder}
\label{sec:Opioid Use Disorder}
Opioid use disorder has a lifetime prevalence of approximately 1--2\% and has been a major driver of the opioid epidemic in North America. The condition results from mu-opioid receptor activation in the ventral tegmental area and nucleus accumbens, producing dopamine release and intense euphoria, with chronic use leading to tolerance, severe withdrawal (agitation, myalgias, lacrimation, rhinorrhea, piloerection, yawning, diarrhea, nausea/vomiting, dilated pupils, muscle spasms, intense craving), and hijacking of the reward system, with respiratory depression being the acute lethal risk. Clinical features include a problematic pattern of opioid use (heroin, prescription opioids) leading to clinically significant impairment, with criteria analogous to AUD. Diagnosis is by clinical interview. Management includes medication-assisted treatment (MAT) as standard of care: methadone (full mu-agonist), buprenorphine (partial mu-agonist), and naltrexone (mu-antagonist), along with psychosocial interventions (CBT, contingency management). Overdose prevention includes naloxone distribution. Prognosis: MAT is associated with reduced mortality; relapse rates after detoxification alone exceed 80\% without ongoing medication.

\subsection{Stimulant Use Disorder (Cocaine, Amphetamines)}
\label{sec:Stimulant Use Disorder (Cocaine, Amphetamines)}
Stimulant use disorder affects approximately 0.5--1\% of the population for cocaine and 0.2--0.4\% for amphetamine-type stimulants. The condition results from cocaine blocking dopamine, norepinephrine, and serotonin transporters, and amphetamines additionally promoting neurotransmitter release from presynaptic vesicles, with chronic use leading to dopamine depletion, neurotoxicity, and extensive structural brain changes. Clinical features include a problematic pattern of stimulant use with criteria analogous to other substance use disorders. Medical complications include acute coronary syndrome, myocardial infarction, hemorrhagic stroke, rhabdomyolysis, seizures, hypertension, and for methamphetamine, severe dental decay (meth mouth), pulmonary hypertension, and psychosis. Withdrawal is characterized by dysphoria, anhedonia, fatigue, hypersomnia, hyperphagia, and intense cravings. Diagnosis is by clinical interview. Management includes behavioral interventions as first-line (CBT, contingency management), with no FDA-approved medications for stimulant use disorder. Prognosis: chronic relapsing course; methamphetamine use disorder has particularly high relapse rates.

%-------------------------------------------------------------
\section{Trauma and Stressor-Related Disorders}
%-------------------------------------------------------------

%-------------------------------------------------------------

%-------------------------------------------------------------

%-------------------------------------------------------------

%-------------------------------------------------------------

%-------------------------------------------------------------

%-------------------------------------------------------------

%-------------------------------------------------------------

\label{sec:Trauma and Stressor-Related Disorders}
%-------------------------------------------------------------%-------------------------------------------------------------

\subsection{Acute Stress Disorder}
\label{sec:Acute Stress Disorder}
Acute stress disorder has a variable prevalence (5--20\% depending on trauma type). The condition describes the development of anxiety, dissociative, and other symptoms within 1 month of a traumatic event, with required symptom clusters similar to PTSD (intrusion, negative mood, dissociation, avoidance, and arousal) but with an emphasis on dissociative symptoms. Clinical features last for 3--30 days after trauma; if symptoms persist beyond 1 month, the diagnosis transitions to PTSD. Diagnosis is by clinical interview. Management includes trauma-focused CBT initiated within the first month. Prognosis is variable; approximately 50--80\% of patients with ASD develop PTSD within 6 months.

\subsection{Adjustment Disorder}
\label{sec:Adjustment Disorder}
Adjustment disorder is one of the most commonly diagnosed psychiatric conditions in medical settings, with prevalence of 5--20\% in outpatient mental health clinics and up to 12\% in hospitalized medical patients. Clinical features include the development of emotional or behavioral symptoms in response to an identifiable stressor within 3 months of onset, with marked distress that is out of proportion to the severity of the stressor, or significant functional impairment, not meeting criteria for another mental disorder. Once the stressor is removed, symptoms typically resolve within 6 months. Diagnosis is by clinical interview. Management is primarily supportive psychotherapy, crisis intervention, and stress management. Prognosis is excellent, with most cases resolving with removal of the stressor.

\subsection{Post-Traumatic Stress Disorder (PTSD)}
\label{sec:Post-Traumatic Stress Disorder (PTSD)}
Post-traumatic stress disorder has a lifetime prevalence of 6--8\% in the general population, with higher rates in combat veterans (10--30\%), sexual assault survivors (30--50\%), and first responders (10--20\%), and is twice as common in women. The condition results from an overly responsive amygdala (hypervigilance, exaggerated fear), a hypoactive medial prefrontal cortex (impaired fear extinction), reduced hippocampal volume, dysregulation of the HPA axis, and heightened sympathetic reactivity, with fear conditioning to trauma-related cues failing to extinguish normally. Clinical features require exposure to actual or threatened death, serious injury, or sexual violence, followed by four symptom clusters for at least 1 month: intrusion symptoms (recurrent, involuntary, intrusive memories, distressing dreams, dissociative reactions, flashbacks), avoidance (persistent avoidance of trauma-related stimuli), negative alterations in cognition and mood, and marked alterations in arousal and reactivity (irritable behavior, hypervigilance, exaggerated startle response, sleep disturbance). Diagnosis is by clinical interview using DSM-5-TR criteria. Management includes trauma-focused psychotherapy as first-line (TF-CBT, prolonged exposure, CPT, EMDR) and pharmacotherapy (SSRIs and SNRIs are FDA-approved). Prognosis is chronic in many cases; about 30--40\% recover fully, 30--40\% have partial improvement, and 20--30\% have persistent symptoms.

